% =============================================================================
% Бүлэг 4: Загвар хөгжүүлэлт ба урсгал
% =============================================================================

\chapter{Загвар хөгжүүлэлт ба урсгал}
\label{ch:model}
\pagecolor{white}

Энэ бүлэгт SQL үүсгэх загвар ба орчуулгын загварыг хөгжүүлсэн дараалал, сургалтын тохиргоо, туршилтын бүлгүүдийг
тайлбарлана.

\section{SQL үүсгэх загвар хөгжүүлэлт ба сургалтын арга зүй}
\label{sec:baseline}

\textbf{Суурь загвар (A бүлэг).} RAG-гүй суурь загварын dev дээрх үр дүнг Хүснэгт~\ref{tab:a-group}-д үзүүлэв. Ижил асуултын
англи хувилбар 43/300 байхад кирилл хувилбар 4/300 байгаа нь \textbf{хэлний зөрүү} гол саад болохыг харуулж байна. Prompt-д
PostgreSQL заахаас (A6) SQLite-аар бичүүлээд sqlglot-оор хөрвүүлэх (A7) нь илүү гарсан тул цаашид A7-г суурь болгов.

\begin{table}[H]
\centering\small
\caption{Суурь загвар (A бүлэг), dev 300, Kaggle T4}
\label{tab:a-group}
\begin{tabularx}{\textwidth}{p{1.2cm}X r}
\toprule
ID & Тохиргоо & Нийт зөв \\
\midrule
A1 & Кирилл асуулт, тайлбаргүй схем, SQLite протокол & 4 (1.3\%) \\
A2 & A1-ийн англи хувилбар & 43 (14.3\%) \\
A3 & A1-ийн латин галиг & 0 \\
A6 & Prompt-д PostgreSQL dialect заах & 6 (2.0\%) \\
A7 & SQLite-аар үүсгээд sqlglot-оор PostgreSQL руу хөрвүүлэх & \textbf{7 (2.3\%)} \\
\bottomrule
\end{tabularx}
\end{table}

\textbf{LoRA тохиргоо.} Rank $r = 32$, $\alpha = 64$, бүх linear давхарга, LR $2\times10^{-4}$, cosine хуваарь, 2 epoch,
loss зөвхөн хариу дээр. A100 дээр bf16-аар сургаж, T4-д ажиллуулахдаа LoRA-г fp32-оор суурь жинтэй нэгтгэж fp16 жин
болгосон (T4 дээр vLLM-ийн LoRA горим ~3 дахин удаан байв). D3 нь Kaggle T4-ийн 16 GB-д багтахын тулд QLoRA ашиглана.

\textbf{Сургалтын урсгал.} \tsql{fetch → translate → synth → dedup → sample → cot → build → sft} алхам бүр өмнөхийн үр
дүнг хадгалж, тасарвал үргэлжилнэ. Kaggle-ийн сесс 12 цагаар хязгаарлагддаг тул хугацааны хязгаарт (\tsql{T2S\_TRAIN\_DEADLINE})
checkpoint хадгалаад зогсож, дараагийн сесс түүнээс үргэлжилдэг болгов (Хүснэгт~\ref{tab:training}).

\begin{table}[H]
\centering\small
\caption{Гүйцэтгэсэн сургалтууд}
\label{tab:training}
\begin{tabularx}{\textwidth}{p{1.1cm}X p{2.2cm} p{3.2cm}}
\toprule
ID & Өгөгдөл & GPU & Төлөв \\
\midrule
C1 & BIRD-ийн монгол орчуулга, rejection sampling & Modal A100 & 380 алхам, ~1 ц 55 мин \\
C2 & C1 + gold-оор чиглүүлсэн CoT & Modal A100 & дууссан \\
C3 & C2 + өөрийн DB-ийн синтетик ба 160 жишээ & Modal A100 & дууссан \\
D2 & B5-ийн RAG prompt (3 756 мөр) & Vast A100 & 470 алхам, ~2.25 ц \\
D3 & Богино prompt, өөрийн DB-ийн жишээ & Kaggle T4×2 & 112/116 алхам, loss 0.555 \\
\bottomrule
\end{tabularx}
\end{table}

\textbf{Туршилтын бүлгүүд.} Бүлэг бүр өмнөх бүлгийн шилдэг тохиргоон дээр \textbf{нэг хүчин зүйлийг} өөрчилнө: A ---
prompt формат, хэл, диалект; B --- RAG-ийн бүрэлдэхүүн (тайлбар B1, утга B2, нэр томьёо B3, жишээ B4, бүгд B5); C, D ---
нэмэлт сургалт; H --- орчуулгын горим ба дүрмийн засвар; E --- засвар (E1), санал хураалт (E2), тодруулга (E3c), шүүлт ба
дүрэм (E4); I --- жишээний сангийн хэмжээ (160 → 579) ба орчуулгын засвар. Нийт 84 нэрлэсэн тохиргоо бий (Хавсралт~\ref{app:configs}).

\begin{table}[H]
\centering\small
\caption{B ба C, D бүлэг, dev 300}
\label{tab:bc-group}
\begin{tabularx}{\textwidth}{p{1.2cm}X r r p{1.2cm}}
\toprule
ID & Тохиргоо & Нийт зөв & Хариулах /224 & GPU \\
\midrule
B1 & A7 + тайлбар, синоним & 11 (3.7\%) & --- & T4 \\
B2 & A7 + утгын хайлт & 14 (4.7\%) & --- & T4 \\
B4 & A7 + 4 жишээ (few-shot) & 22 (7.3\%) & --- & T4 \\
B5 & B1--B4 бүгд (RAG) & \textbf{43 (14.3\%)} & --- & T4 \\
C1 & LoRA SFT, rejection sampling, RAG-гүй & 9 (3.0\%) & 7 & A100 \\
C3 & C2 + өөрийн DB-ийн өгөгдөл & 30 (10.0\%) & 26 & A100 \\
D2 & RAG prompt-оор сургасан + RAG & 47 (15.7\%) & 37 & A100 \\
\bottomrule
\end{tabularx}
\end{table}

\textbf{Бүрэлдэхүүн тус бүрийн нөлөө.} Хүснэгт~\ref{tab:ablation}-т бүрэлдэхүүн бүрийг нэмэхэд dev-ийн нийт зөв хэрхэн
өөрчлөгдсөнийг үзүүлэв. Мэдлэг хайх ба орчуулга хамгийн их, засвар (E1) ба өөрийн DB-ийн жишээгүй нэмэлт сургалт хамгийн
бага нэмэр өгсөн.

\begin{table}[H]
\centering\small
\caption{Бүрэлдэхүүн тус бүрийн нөлөө (dev 300, нийт зөв, \%)}
\label{tab:ablation}
\label{sec:ablation}
\begin{tabularx}{\textwidth}{p{2.4cm}X r r}
\toprule
Шилжилт & Нэмсэн бүрэлдэхүүн & Өмнө → дараа & Δ \\
\midrule
A7 → B5 & Мэдлэг хайх (RAG) & 2.3 → 14.3 & +12.0 \\
B5 → D2 & LoRA нэмэлт сургалт (өөрийн DB-ийн жишээгүй) & 14.3 → 15.7 & +1.4 \\
D2 → D2 tr & Англи орчуулга & 15.7 → 27.0 & +11.3 \\
D2 → H1 & Дүрмийн засвар, татгалзал & 15.7 → 24.7 & +9.0 \\
H1 → H4 & Орчуулга: загварт англи, хайлт монголоор & 24.7 → 36.7 & +12.0 \\
H4 → E1 & Алдааны мессежээр засах & 36.7 → 37.0 & +0.3 \\
E1 → E2 & 8 хувилбар, санал хураах & 37.0 → 41.0 & +4.0 \\
E2 → E3c & Итгэл багатай бол тодруулах & 41.0 → 45.7 & +4.7 \\
E2 → E4 & Хувилбарын шүүлт + dev-ээс гаргасан дүрэм & 41.0 → 53.0 & +12.0$^\ast$ \\
\bottomrule
\end{tabularx}
\par\smallskip\footnotesize $^\ast$ E4-ийн дүрмийг dev-ийг уншсаны дараа бичсэн тул dev-ийн тоо хэт өөдрөг; test дээр
+6.5 пункт (§\ref{sec:test-results}).
\end{table}

\section{Орчуулгын загвар хөгжүүлэлт ба сургалтын арга зүй}
\label{sec:mt-train}

\textbf{Zero-shot орчуулга.} Эхний шатанд Qwen3.5-9B-д орчуулах заавар өгч ашиглав. DB-д хадгалагдсан утгыг
(«Улаанбаатар», «Мобайл апп») орчуулахаас өмнө \tsql{[V1]} маягийн маскаар сольж, дараа нь сэргээнэ. Хоёр дахь хувилбар
(T2) маскыг зөвхөн DB-д үнэхээр байгаа нэрээр шүүж, мэдлэгийн сангаас нэр томьёоны орчуулгыг (≤6 мөр) зааварт нэмдэг;
энэ нь 500 асуултын маскыг 122-оос 37 болгож бууруулсан. Хайлтыг монгол асуултаар, SQL загварт өгөх асуултыг англиар
хийх (H4) хослол хамгийн сайн гарсан.

\textbf{Нэмэлт сургалт (M1).} Орчуулгаас үүдэх алдааг багасгахын тулд MADLAD-400-3B-MT-г MT v2 өгөгдлөөр (8 601 мөр)
QLoRA-аар нэмж сургана (Хүснэгт~\ref{tab:mt-hparams}). Сургалтын өмнө нэмж сургаагүй загварын dev орчуулгыг (M0)
хэмжиж, сургалтын дараа val хосоор chrF++, dev асуултаар нэр томьёо ба Text2SQL-ийн EX-ийг харьцуулна.

\begin{table}[H]
\centering\small
\caption{MADLAD-400-3B-MT-ийн QLoRA сургалтын гипер-параметр}
\label{tab:mt-hparams}
\begin{tabularx}{\textwidth}{p{4.6cm}X}
\toprule
Параметр & Утга \\
\midrule
Квантчлал & 4 бит NF4, double quant; тооцоолол fp32 (T5-ийн fp16 overflow-оос сэргийлж; T4-д bf16 алга) \\
LoRA & $r = 32$, $\alpha = 64$, dropout 0.05; q, k, v, o, wi\_0, wi\_1, wo \\
Оновчлол & LR $3\times10^{-4}$, cosine, warmup 3\%, 2 epoch (270 алхам) \\
Batch & 4 × градиент хуримтлал 8 (= 32), дээд урт 256 токен (batch 16 нь T4-д санах ой хүрэлцээгүй) \\
Сургагдах параметр & 94.4 сая (нийтийн 3.1\%) \\
Үнэлгээ & val 287 хос (кирилл, латин), beam 4, chrF++; суурь нь сургалттай ижил NF4 \\
\bottomrule
\end{tabularx}
\end{table}

Сургалтыг Kaggle T4 дээр 2026-10-26-нд хийсэн (Version 16, 3.7 цаг; loss 43.3 → 7.9). Үүнээс өмнөх оролдлогууд дараах
шалтгаанаар зогссон: кирилл үсгийг байтаар таслах ба хоосон объект дуудах скриптийн алдаа (Version 11--12), transformers~5
MADLAD-ийн embedding-ийг холбохгүй ачаалснаас хог гаралт гарах (Version 13--14; туршилтын өгүүлбэрээр зөв embedding-ийг
автоматаар сонгох засвар хийсэн), хоёр GPU-д хуваасан beam search-ийн CUDA алдаа ба санах ой хүрэлцээгүй (Version 15--17).
Нэмж сургасан загварыг fp32 суурьтай ачаалахад T4-ийн 15~GB-д багтаагүй тул M1-ийн орчуулгыг сургалттай ижил NF4 суурь дээр
хийсэн; M0 нь fp32 суурьтай.

\section{Бүлгийн дүгнэлт}
\label{sec:ch5-summary}

Arctic-R1-7B-ийн сургалтын протоколыг мөрдөж, SQLite → PostgreSQL хөрвүүлэлттэй суурийг тогтоов. Хязгаарлагдмал GPU дээр
LoRA ба QLoRA сургалтыг тасалдлаас сэргэх боломжтойгоор хэрэгжүүлэв. Бүрэлдэхүүнээр нь задлахад RAG ба орчуулга тус бүр
11--12 пункт нэмсэн бол өөрийн DB-ийн жишээг хассан нэмэлт сургалт 1.4 пункт л нэмсэн. Орчуулгын загварыг санхүүгийн
хэллэгээр нэмж сургах ажил эхэлсэн.
